\section{Methodological Development and Decisions}
\label{sec:methodologicaldevelopment}

The experiments were developed in stages to separate graph recovery,
prediction, structural intervention effects and behavioral measurement. We
report this sequence because a successful check at one stage did not establish
the claim tested at the next. Trials~1--5 below contain completed evidence;
Trial~6 adds a completed exploratory compression diagnostic and identifies the
remaining proposed encoder studies. Changes motivated by inspected development
data are not counted as independent confirmation, and the original failed
criteria remain part of the record.

\subsection{Trial 1: graph recovery did not establish useful endpoint prediction}

We first checked that the imported conventional and fast c-GC variants could
recover simple generating graphs before interpreting downstream failures
\citep{adedayo2025cgc,adedayo2026fastcgc}. All 24 method/scenario/seed runs in the
chain and independent-null technical panel recovered the expected support
(Section~\ref{sec:fourcslsanity}). This established a working native-estimation
path under those favorable controls; it did not validate a behavioral manifold
or an intervention readout.

We then tested whether graph features added predictive value in a shared,
training-owned manifold-transition comparison. The framework pilot completed
48 method/scenario/seed runs: four methods, four scenarios and three process
seeds. None of its 16 method/scenario contrasts supplied multiplicity-adjusted
evidence of improvement; all Holm-adjusted $p$-values were one. With only three
independent seeds, this pilot had limited inferential resolution. We therefore
retained the narrow technical success and required a separate simulator-grounded
test of the claimed driver quantity, rather than treating graph recovery or
geometric correspondence as evidence of behavioral causation.

\subsection{Trial 2: feature omission failed the structural driver test}

The subsequent driver study completed all 560 planned method/recording runs,
but none of the four methods met the frozen recovery and null criteria. Its
candidate score removed outgoing graph-feature blocks and refitted a predictor
while retaining the remaining activity and manifold-history information. The
gold standard instead disconnected outgoing influence in the generating neural
mechanism. Conditional predictive importance and the effect of a structural
intervention answer different questions
\citep{williamson2023importance,pearl2010causalinference}.

We consequently revised the intervention readout, not the completed study's
acceptance rules. The estimand mismatch is a reason to test an explicit neural
mechanism; these experiments do not establish it as the sole cause of every
failure in the earlier latent benchmark. Its observation process, endpoint and
selection problem also differ from the favorable control introduced next. The
560-run negative result is therefore retained under its original protocol
(Section~\ref{sec:sharedframeworkvalidation}).

\subsection{Trial 3: an explicit mechanism worked when discovery history was adequate}

We implemented the fully observed additive-tanh mechanism in
Section~\ref{sec:critical}, retaining target self-history and estimating
coefficients and innovation scales from training trials. Outgoing disconnection
changes cross-neuron update terms while preserving the sources' own evolution.
Known-support and empty-support controls separate mechanism learning from
support-discovery error. Regularization and residual-noise estimation use inner
whole-trial validation, followed by refitting on the outer training trials
\citep{hoerl1970ridge,cawley2010overfitting}.

An initial paired check held the recording, full mechanism-fitting data,
evaluation origins and rollout noise fixed while changing native discovery
history from 256 to 1024 samples. At development seed~20, all four native
methods failed the short-history driver ranking but recovered the generating
support and ranking with the longer history. Long-history effect MAE was
0.001831, compared with 0.140055 for empty support
(Section~\ref{sec:mechanismdevelopment}). This localized an important
finite-history limitation in this control. It did not justify lengthening
histories across arbitrary biological regime changes. Because this comparison
had only one process seed, its favorable result required fresh replication
and a separate null assessment.

\subsection{Trial 4: fresh replication supported effect estimation but retained an eligibility failure}

The next study separated 19 calibration-null seeds, 20 fresh positive-test
seeds and 60 untouched null-test seeds. Its primary history was 1024 samples;
the same positive-test recordings also received a paired 256-sample discovery
comparison. All 119 cells and 714 method/control runs completed, including
4284 native window fits. Calibration was completed before test evaluation,
retaining the minimum probability-effect cutoff of 0.10.

At the primary history, all four native methods attained mean effect MAE
$0.002087\pm0.000805$ across the 20 positive-test seeds; the uncertainty shown
here is the sample SD of seed means. Mean MAE was 97.55--98.36\% lower than
with the paired shorter discovery history. Each method had zero false-event
seeds among the 60 null tests. The per-method one-sided exact 95\% upper bound
was 0.048703, meeting the declared 0.05 null gate
\citep{clopper1934confidence}. This bound concerns one method's whole-seed
false-nomination event; it is not simultaneous control across the four variants
or general neuron-level false-discovery control.

The complete frozen acceptance decision nevertheless remained false. Only
19 of the 20 positive seeds satisfied individual-ranking eligibility. In
seed~20012, one source condition had no individual neuron above the fixed
gold-effect threshold, although the predeclared driver-pair intervention had
effect 0.230469 and was correctly estimated and nominated by all four methods.
Known support had the same individual-ranking ineligibility. We did not lower
the gold threshold, exclude that condition or replace the complete-set gate
with the perfect eligible-subset description. The result supports accurate
outgoing-disconnection effects in this model-matched control and identifies a
need to distinguish individual and joint causal units prospectively. It does
not validate biological silencing or a learned behavioral endpoint.

\subsection{Trial 5: the learned chart changed the behavioral question}

We next replaced the directly observed endpoint with training-frozen learned
measurements while keeping the structural model in full neural space. This
separates the measurement question from an additional assumption that projected
coordinates possess autonomous dynamics. PCA and geometric regions describe
neural variation; their behavioral meaning must be tested rather than assumed
\citep{jolliffe2016pca}. The independent reference was the generator's
prespecified behavioral coordinate and event, not a principal component or a
cluster selected after evaluation.

Five development seeds (40000--40004) each supplied driver and endpoint-null
cases. All 10 cases and 70 dynamics runs completed, with 360 native window fits.
Each neural model and its paired factual/intervened futures were reused across
the seven endpoint arms. Encoders and geometric regions used inner-training
trials; behavioral decoder regularization used inner-validation trials. The
encoder was then frozen while the decoder was refitted on outer training data.
All comparisons retained the same held-out factual origins and included source
classification errors in the transition event. Re-executing these five seeds
reproduced the earlier development errors and did not add independent evidence.

We distinguished three quantities: fitted-versus-simulator effect error through
the same chart; simulator-chart versus independently specified behavioral
effect error; and the final fitted-chart versus true-behavior effect error.
Table~\ref{tab:developmentendpointdiagnosis} shows why the first alone is
insufficient. All 40 native consensus supports matched known support, and all
four native methods had the same errors as the known-support fit in this panel.
The true simulator itself had large PCA semantic error, so more accurate
mechanism fitting could not remove the measurement discrepancy.

\begin{table}[htbp]
\centering\small
\caption{Driver-scenario endpoint diagnosis, five development seeds. Entries
are means of seed-level absolute probability-effect errors over identical
origins/candidates; lower is better. Within-chart error compares fitted and
simulator effects through the frozen chart. Semantic error compares the
simulator's chart and behavioral effects. Behavioral error compares fitted chart
effects with true behavioral effects. Known support and all four native methods
have identical entries here; they are paired conditions, not extra replicates.}
\label{tab:developmentendpointdiagnosis}
\begin{tabular}{@{}lrrr@{}}
\toprule
Endpoint & Within chart & Semantic & Behavioral \\
\midrule
PCA geometric regions & 0.002018 & 0.130737 & 0.131144 \\
PCA + behavioral decoder & 0.002181 & 0.138428 & 0.138818 \\
Full-neural behavioral decoder & 0.002067 & 0.000553 & 0.001969 \\
\bottomrule
\end{tabular}
\end{table}

The PCA-plus-decoder arm reached approximately 90\% natural future-target
classification accuracy but only 47\% factual-origin accuracy in the driver
case. Its final behavioral effect error remained 0.138818. Behavioral
supervision added after this unsupervised compression therefore did not repair
the endpoint. Constant and source-proxy charts could have zero within-chart
effect error while missing true behavioral effects; the constant chart also
failed occupancy. A neural-nuisance chart could respond to intervention under
the behavioral null. These controls expose measurement failure that chart
agreement alone would conceal. No endpoint-calibrated nomination or nonlinear
manifold claim was established by this five-seed bridge.

\subsection{Trial 6: behavior-prioritized compression as a capacity diagnostic}

We next asked whether the failed PCA endpoint implied that a compact outcome
measurement was impossible. For the saved affine encoder $q(X)=AX+b$ and frozen
full-neural linear classifier with weights $w$ and intercept $c$, we formed
\begin{equation}
 z(X)=w^{\top}(AX+b)+c,\qquad
 \widehat B(X)=\mathbf 1\{z(X)\geq0\}.
\end{equation}
This scalar supervised coordinate is an algebraic re-expression of the existing
decoder. Its future target thresholds the mean affine score, not a mean of
sigmoid probabilities; linearity preserves the previous mean-code classifier
and the same zero-tie convention. The exploratory probe reused the same five
inspected development seeds, 10 cases and 70 saved dynamics parameter sets,
without new native discovery or mechanism fits.

The scalar score changed none of the parent decoder's labels on 276480
naturally generated neural observations or its 1935360 paired-rollout binary
events. Its behavioral MAE was 0.001969 in the driver case and 0.000073 in the
null, matching the full-neural decoder. These counts audit the implementation;
they do not add independent observations or increase the five-seed sample size.
Any binary linear classifier admits such a scalar decision score. Its existence
therefore demonstrates capacity to preserve this particular readout, not an
independently improved learning algorithm, intrinsic behavioral dimension,
manifold geometry or closed latent dynamics. Computing the score still requires
the full observed neural vector, and the structural simulator remains in that
full space.

We also disclosed the complete one- through five-component PCA-plus-decoder
sweep. Driver behavioral MAEs were respectively 0.146908, 0.138818, 0.138875,
0.138818 and 0.018994. Four components retained a mean 90.45\% of standardized
neural training variance in the driver cases, yet their effect error remained
large. More retained variance alone did not ensure the relevant measurement;
the inspected sweep is not a test-selected winning dimension or fresh
confirmation. The score factorization and dimension sweep are preserved in
\path{synthetic/results/behavior_prioritized_probe_2026-11-18/}.

The resulting next study remains proposed: learn the encoder using a
prespecified, distributed behavioral readout. We considered a supervised linear
or partial-least-squares baseline; preferential subspace identification for
behaviorally relevant temporal dynamics \citep{sani2021psid}; CEBRA's
behavior-informed or hybrid representation \citep{schneider2023cebra}; and joint
behavioral/dynamical BunDLe-Net or a supervised autoencoder
\citep{kumar2023bundleV5,le2018supervisedautoencoders}. These address different
objectives and do not automatically validate outgoing-mechanism interventions.
Temporal encoders require genuine causal prehistory rather than future
borrowing. A new behavioral reference also requires an oracle-checked null:
the present null removes inputs only to the final coordinate and is not
automatically null for a readout containing the responsive bystander. Nonlinear
alternatives have not been run. Frozen-map semantic checks, followed by new
calibration and test pools, remain necessary. Graph similarity is a separate
grouping option; correct support already coexisted with the failed PCA endpoint
in this control.
